\documentclass{report}
\input{preamble}
\usepackage[style=alphabetic, backend=biber]{biblatex}
\usepackage{tocloft}
\addbibresource{references.bib}
\setlength{\cftbeforesecskip}{0.3em}
\setcounter{tocdepth}{2}

\title{Reinforcement Learning}
\author{Hyoung Joon, Kim}
\date{\today}

\begin{document}
\maketitle
\tableofcontents
\newpage
\chapter{Introduction}
        This document is a set of notes of algorithms of reinforcement learning and theories of reinforcement learning.
        Following books were mainly used for studying algorithms and theories.
        \begin{itemize}
            \item \cite{Sutton2020};
            \item \cite{Graesser2020};
            \item \cite{Zhao2025};
        \end{itemize}
\part{Common Theories}
    \chapter{Basic Concepts}
    \section{Introduction}
In this section, we discuss common concepts that are used frequently in reinforcement learning.

\section{State, Action, Policy and Reward}
\subsection{State}
    The \textbf{state} is the status of the agent with respect to the environment.
    The set of all possible states, the \textbf{state space}, is denoted as $\mathcal{S}$.
\subsection{Action}
    The \textbf{action} is what agent can do, when a state is given. By choosing the action and
    taking it, the agent moves to the next state, which is called \textbf{state transition}. The set of all possible actions in state $s$,
    \textbf{action space} for state $s$, is denoted as $\mathcal{A}(s)$.
    \subsubsection{Transition Probability}
        Note that the transition can be stochastic. That is, the same choice of action in the same state
        may not lead to equivalent next state. We denote this probability with conditional porbability as following:
        \begin{equation}
            p(s'|s,a)
        \end{equation}
        where $s',s \in \mathcal{S}$ and $a \in \mathcal{A}(s)$.
\subsection{Policy}
    The \textbf{policy} tells the agent which actions to take at every state. The policy can be deterministic or stochastic.
    We denote the probability of taking action $a$ under state $s$ following policy $\pi$ as $\pi(a|s)$.
\subsection{Reward}
    \textbf{Reward} is one of the most unique concept of reinforcement learning. After the execution of action, the agent
    receives a reward, denoted as $r$. The agent's goal is to maximize the total reward that the agent receives.
    Due to this goal, the reward can encourage or discourage the agent's
    choice of selecting an action under certain state. Again, the reward can be stochastic, too.

    One question that follows naturally is: Can't we just pick the actions that gives the most reward at every state? If this was correct,
    the reinforcement learning may not have existed. Since the reward is just an immediate information of the act, choosing only the action with
    highest reward is not the correct answer, most of the time. There can be a better choice in the long run.

\section{Trajectories, Returns and Episodes}
    \subsection{Trajectory}
        The \textbf{trajectory} is the footsteps of the agent. More formally, suppose that the agent is on the initial state $s_0$ and took action $a_0$.
        The agent receives reward $r_1$ from environment and moves to state $s_1$. Again, the agent take an action $a_1$, receive reward $r_1$ and move to state $s_2$.
        Repeating this, we get the trajectory of an agent:
        \begin{equation}
            s_0,a_0,r_1,s_1,a_1,r_2,\dots,s_{T-1},a_{T-1},r{T}. \label{eq:trajectory}
        \end{equation}
    \subsection{Returns}
        Given the trajectory $\tau = \{s_0,a_0,r_1,s_1,a_1,r_2,\dots,s_{T-1},a_{T-1},r{T}\}$, we can calculate the \textbf{return} $G_t$, of the trajectory.
        The naive definition of the return, called \textbf{total rewards} is the sum
        \begin{equation}
            G_t = \sum_{k=t+1}^{T}r_k. \label{eq:total-reward}
        \end{equation}
        The point is that we do not know if the terminal time $T$ is finite or not. If $T$ is not finite, than the cummulative reward \ref{eq:total-reward}
        diverges. For this, we introduce the \textbf{discounted return}, defined as
        \begin{equation}
            G_t = \sum_{k=t+1}^{\infty}\gamma^{k - t - 1}r_k
        \end{equation}
        where $\gamma \in [0,1]$ is called \textbf{discount rate}.
        This works for for finite trajectory too, if we let $r_k = 0$ for $k > T$.
    \subsection{Episodes}
        The \textbf{episode} is the finite trajectory itself, like \ref{eq:trajectory}. The episode may stop at the states called \textbf{terminal states}.
        The tasks with episodes are called \textbf{episodic tasks}. Unlike episodic tasks, there are \textbf{continuing tasks}, which does not ends. These two can
        be expressed in unified mathematical manner. For this, in the episodic task, we must consider the agent's behavior after the terminal state:
        to stay in that state forever, or to take action like normal states. In this note, we follow the latter, like \cite{Zhao2025}.

\section{Markov Decision Process}
    \subsection{Definition}
        The \textbf{Markov decision process} (MDP) is the general framework that describes the stochastic dynamic systems. It is defined as following:
        \begin{defn}{Markov Decision Process}{MDP}
            The \textbf{Markov decision process} is the 5-tuple $(\mathcal{S},\mathcal{A},p,\mathcal{R},\gamma)$, where each elements are defined as following:
            \begin{itemize}
                \item $\mathcal{S}$ is the state space, $\mathcal{A}(s)$ is the action space associated with state $s$, and $\mathcal{R}(s,a)$ is a set of rewards, associated with each state-action pair $(s,a)$.
                \item $p(s',r|s,a)$ is the dynamics of the model, which is as probability of moving to state $s'$ from $s$ by taking action $a$, and obtaining reward $r$.
                Note that the state transition probability and the reward probability can be decribed with dynamics. See \ref{rem:dynamics}
                \item $\pi(s|a)$ is the policy, a probability of agent choosing action $a$ in state $s$.
                \item $\gamma$ is the discount rate.
            \end{itemize}
            with the dynamics satisfying the \emph{Markov property}:
            \begin{align}
                p(s_{t+1}|s_t,a_t,s_{t-1},a_{t-1},\dots,s_0,a_0) &= p(s_{t+1}|s_{t},a_{t}) \\
                p(r_{t+1}|s_t,a_t,s_{t-1},a_{t-1},\dots,s_0,a_0) &= p(r_{t+1}|s_{t},a_{t})
            \end{align}
            where $t$ is the current time step.
        \end{defn}
        \begin{rem}{Transition and reward probability as dynamics}{dynamics}
            The 4-tuple dynamics can be used to describe transition probability and reward probability as following:
            \begin{align}
                p(s'|s,a) &= \sum_{r \in \mathcal{R}(s,a)}p(s',r|s,a) \\
                p(r|s,a) &= \sum_{s' \in \mathcal{S}}p(s',r|s,a)
            \end{align}
        \end{rem}
        The dynamics can be either stationary or non-stationary. If the dynamics changes over time, it is non-stationary; else, it is called stationary.

    \chapter{Bellman Operator}
    \section{Introduction}
    This section introduces the Bellman equation and the Bellman optimality equation,
    using the Bellman operator as a main framework. This part was influenced from \cite{littman1996}.
    Note that we always assume that the state space and action space is finite. Although there are algorithms that can be applied to continuous states or actions (\ref{alg:DQN}),
    for simplicity, we assume finiteness. (and actually, no convergence is guaranteed for DQN.)
\section{The State Value Function and Action Value Function}
    Before diving into the Bellman equation, we introduce an important concept called the \textbf{value of a state} and the \textbf{value of an action}.
    \begin{defn}{State Value}{state-value}
        Let the agent follows the policy $\pi$. The \textbf{state value} $v_\pi(s)$ of state $s$ is defined as the expected return:
        \begin{equation}
            v_\pi(s) = \E[G_t|S_t=s] \label{eq:state-value}
        \end{equation}
        where $G_t$ is the discounted return $G_t=\sum_{k=t+1}^{\infty}\gamma^{k-t-1}R_{k}$.
        The function $v_\pi$ itself is called \textbf{state-value function}.
    \end{defn}
    State values are used to evaluate the policy $\pi$. Policies that generate greater state values are better. Then, how can we
    obtain the state values? The answer is to solve the Bellman equation.
    \begin{defn}{Action Value}{action-value}
        Let $\pi$ be a given policy. Then, the \textbf{action value} of a state-action pair is defined as
        \begin{equation}
            q_\pi(s,a) = \E[G_t|S_t=s,A_t=a].
        \end{equation}
    \end{defn}
    \begin{rem}{Properties of Action Value}{ActionValProp}
        Action value $q_\pi(s,a)$ has following properties.
        \begin{itemize}
            \item $v_\pi=\sum_{a \in \mathcal{A}(s)} \pi(a|s)q_\pi(s,a)$
            \item $q_\pi(s,a) = \sum_{r \in \mathcal{R}(s,a)}p(r|s,a)r + \gamma \sum_{s' \in \mathcal{S}}p(s'|s,a)v_\pi(s')$
        \end{itemize}
    \end{rem}
    \begin{proof}
        \begin{itemize}
            \item Use the law of total expectation \ref{eq:cond-LTE}.
            \item Use the result above and \ref{thm:BellmanEq-SV}.
        \end{itemize}
    \end{proof}
    Actually, the ~\ref{rem:ActionValProp} is the \textbf{Bellman equation} of state values.
\section{Generalized Bellman Equation}
    \begin{defn}{Generalized Bellman Equation}{GeneralizedBellmanEq}
        Let $v(s)$ be a state value function. Then, the \textbf{generalized Bellman equation} is defined as following:
        \begin{equation}
            v(s) = \bigotimes_a \Big(r(s, a) + \gamma \bigoplus_{s'} v(s') \Big) 
        \end{equation}
        where $\bigotimes$ defines how an optimal agent should choose actions, and $\bigoplus$ defines how the value of the next state should be used in assigning value to the current state.

        From here, we can define the \textbf{generalized Bellman operator} $T:\mathbb{R}^\mathcal{S} \to \mathbb{R}^\mathcal{S}$ as following:
        \begin{equation}
            [Tv](s) = \bigotimes_a \Big( r(s,a) + \gamma \bigoplus_{s'}v(s') \Big)
        \end{equation}

        The \textbf{myopic policy} $\pi$ with respect to a value function $v$ is any (stochastic) policy such that $T^\pi v = T v$ where
        \begin{equation}
            [T^\pi v](s) = \sum_a \pi(s|a) \Big( r(s, a) + \gamma \bigoplus_{s'}v(s') \Big)
        \end{equation}
    \end{defn}
    For a myopic policy to exist, we require that the operator $\bigotimes$ satisfy the following property: for all functions $f$ and states $s$,
    $\min_a f(s,a) \leq \bigotimes_a f(s,a) \leq \max_a f(s,a)$ (\cite{littman1996}).
    And we will only use the operator $\Big[\bigotimes_{s'} v\Big](s) = \sum_{s' \in \mathcal{S}} p(s'|s,a) v(s')$, since we only cope with
    MDPs.
\section{Bellman Equation}
    The value function with respect to a policy $\pi$ satisfies the following Bellman equation.
    The Bellman equation is simply a set of linear equations that describe the relationships of each state values.
    Bellman equation can be obtained by using the recursive property of return:
    \begin{align}
        G_t &= R_{t+1} + \gamma R_{t+2} + \gamma^2 R_{t+3} + \gamma^3 R_{t+4} + \cdots \\
            &= R_{t+1} + \gamma (R_{t+2} + \gamma R_{t+3} + \gamma^2 R_{t+4}) + \cdots \\
            &= R_{t+1} + \gamma G_{t+1}.
    \end{align}
    Substituting equation \ref{eq:state-value} with the recursive realationship, we obtain
    \begin{align}
        v_\pi(s) &= \E[G_t|S_t=s] \\
                    &= \E[R_{t+1} + \gamma G_{t+1}|S_t=s] \\
                    &= \E[R_{t+1}|S_t=s] + \gamma \E[G_{t+1}|S_t=s]
    \end{align} 
    Using the law of total expectation \ref{eq:cond-LTE}, we get
    \begin{align}
        \E[R_{t+1}|S_t=s] &= \sum_{a \in \mathcal{A}(s)}\pi(a|s)\E[R_{t+1}|S_t=s,A_t=a] \\
                            &= \sum_{a \in \mathcal{A}(s)}\pi(a|s)\sum_{r \in \mathcal{R}(s,a)}p(r|s,a)r
    \end{align}
    and
    \begin{align}
        \E[G_{t+1}|S_t=s] &= \sum_{s' \in \mathcal{S}}\E[G_{t+1}|S_{t+1}=s',S_t=s]p(s'|s) \\
                            &= \sum_{s' \in \mathcal{S}}\E[G_{t+1}|S_{t+1}=s']p(s'|s) \quad \text{(Markov property)} \\
                            &= \sum_{s' \in \mathcal{S}}v_\pi(s')p(s'|s) \\
                            &= \sum_{s' \in \mathcal{S}}v_\pi(s')\sum_{a \in \mathcal{A}(s)}p(s'|s,a).
    \end{align}
    The final result is given in the following box:
    \begin{thm}{Bellman Equation for State-Value Function}{BellmanEq-SV}
        Let $\pi$ be a policy, and let $v_\pi$ be a state-value function. Then, the following Bellman equation is satisfied:
        \begin{equation}
            v_\pi(s) = \sum_{a \in \mathcal{A}(s)}\pi(a|s) \Big(r(s,a) + \gamma \sum_{s' \in \mathcal{S}}p(s'|s,a)v_\pi(s')\Big)
        \end{equation}
    \end{thm}

    And one can see the operator $\Big[\bigotimes_{s'} v\Big](s) = \sum_{s' \in \mathcal{S}} p(s'|s,a) v(s')$.

    \subsection{Matrix-Vector form of Bellman Equation}
        We can reformulate Bellman equation in terms of matrix-vector form. For this, we rewrite the Bellman equation given in Theorem \ref{thm:BellmanEq-SV} as
        \begin{equation}
            v_\pi(s) = r_\pi(s) + \gamma\sum_{s' \in \mathcal{S}}p_\pi(s'|s)v_\pi(s'),
        \end{equation}
        where
        \begin{align}
            r_\pi(s) &= \sum_{a \in \mathcal{A}(s)}\pi(a|s)\sum_{r \in \mathcal{R}(s,a)}p(r|s,a)r \\
            p_\pi(s'|s) &= \sum_{a \in \mathcal{A}(s)}\pi(a|s)p(s'|s,a).
        \end{align}
        The $r_\pi$ and $p_\pi$ denotes the mean of immediate rewards and probability of transitioning from $s$ to $s'$ under policy $\pi$, respectively.
        Suppose that we indexed all states as $s_i$ with $i = 1, \dots, n$, where $n=|\mathcal{S}|$.
        Let
        \begin{equation}
            v_\pi = \begin{bmatrix}
                    v_\pi(s_1) \\
                    \vdots \\
                    v_\pi(s_n)
                    \end{bmatrix},
            r_\pi = \begin{bmatrix}
                    r_\pi(s_1) \\
                    \vdots \\
                    r_\pi(s_n)
                    \end{bmatrix},
            P_\pi = \begin{bmatrix}
                    p_\pi(s_1|s_1) & p_\pi(s_2|s_1) & \dots & p_\pi(s_n|s_1) \\
                    p_\pi(s_1|s_2) & p_\pi(s_2|s_2) & \dots & p_\pi(s_n|s_2) \\
                    \vdots & \vdots &  & \vdots \\
                    p_\pi(s_1|s_n) & p_\pi(s_2|s_n) & \dots & p_\pi(s_n|s_n)
                    \end{bmatrix}. \label{eq:VecForm}
        \end{equation}
        Then, Theorem \ref{thm:BellmanEq-SV} can be written in the following form:
        \begin{thm}{Matrix-Vector form of Bellman Equation}{MVFormBellmanEq}
        Let $v_\pi$, $r_\pi$ and $P_\pi$ be given as \ref{eq:VecForm}. Then, following is satisfied:
            \begin{equation}
                v_\pi = r_\pi + \gamma P_\pi v_\pi
            \end{equation}
        \end{thm}
        \begin{rem}{Properties of Matrix $P_\pi$}{ProbMatProp}
            Note that the matrix $P_\pi$ given in equation \ref{eq:VecForm} satisfies following properties:
            \begin{itemize}
                \item All the elements are nonnegative;
                \item The sum of the values of every row equals 1. Matrix with this property is called a \emph{stochastic matrix}.
            \end{itemize}
        \end{rem}

\section{Bellman Optimality Equation}
    \subsection{Optimal Policy}
        The reinforcement learning algorithms are all designed to find an optimal policy for
        given enviroment. Then, what is an 'optimal' policy? The definition is given below.
        \begin{defn}{Optimal Policy and Optimal State Value}{optPol}
            A policy $\pi^*$ is optimal if and only if $v_\pi^*(s) \geq v_\pi(s)$ for all $s \in \mathcal{S}$
            and for any other policy $\pi$. The state values of $\pi^*$ are the optimal state values.
        \end{defn}
        \begin{defn}{Better Policy}{betterPol}
            A policy $\pi_1$ is \textbf{better} than a policy $\pi_2$ if and only if
            $v_{\pi_1}(s) \geq v_{\pi_2}(s)$ for all $s \in \mathcal{S}$.
        \end{defn}
        The questions that naturally arise are:
        \begin{itemize}
            \item Existence: Does the optimal policy exists?
            \item Uniqueness: Is the optimal policy unique?
            \item Stochasticity: Is the optimal policy stochastic or deterministic?
            \item Algorithm: How do we obtain the optimal policy and the optimal state values?
        \end{itemize}
    \subsection{Bellman Optimality Equation}
        The optimal policy and optimal state values can be analyised via the \textbf{Bellman optimality equation} (BOE).
        The optimal state values satisfies following Bellman optimality equation. This claim in fact has two logically independent statements:
        \begin{enumerate}
            \item The BOE, viewed purely as a fixed-point equation, has a unique solution (Section \ref{sec:SolveBellmanOptEq});
            \item The unique solution of BOE concides with the optimal stat-value function $v^*$, and the derived greedy policy is optimal (Section \ref{sec:OptPolOptStateValBOE})
        \end{enumerate}
        Note that the Section \ref{sec:SolveBellmanOptEq} does not assume that the state-value function is optimal.
        \begin{defn}{Bellman Optimality Equation}{BellOptEq}
            Let $\Pi$ be a set of possible policies and let $\gamma \in (0,1)$ be a discount rate. Then, the Bellman Optimality Equation is defined as
            \begin{align}
                v(s) &= \max_{a \in \mathcal{A}(s)}\left( r(s,a) + \gamma \sum_{s' \in \mathcal{S}}p(s'|s,a)v(s') \right) \\
                    &= \max_{a \in \mathcal{A}(s)} q(s,a)
            \end{align}
            where $v(s),v(s')$ are unkown variables to be solved.

            The Bellman optimality operator $T^*:\mathbb{R}^\mathcal{S} \to \mathbb{R}^\mathcal{S}$, uses $\bigotimes_a = \max_a$ as its action selection strategy.
            Specifically, the Bellman optimality operator is defined as following:
            \begin{equation}
                [T^* v](s) = \max_{a \in \mathcal{A}(s)}\left( r(s, a) + \gamma \sum_{s' \in \mathcal{S}}p(s'|s,a)v(s') \right)
            \end{equation}
        \end{defn}
        
\section{Solution of The Bellman Equations}
    \subsection{Contraction Mapping Theorem}
        We propose a \textbf{contraction mapping theorem}, the main tool for solving the Bellman equations.
        The contraction mapping theorem is a powerful tool for analyzing general nonlinear equations.
        \begin{defn}{Contraction Mapping}{ContractMap}
            Let $f : \mathbb{R}^d \to \mathbb{R}^d$. Then, the function $f$ is a \textbf{contraction mapping}
            if there exists $\gamma \in (0,1)$ such that
            \begin{equation}
                \|f(x_1) - f(x_2)\| \leq \gamma \|x_1-x_2\|
            \end{equation}
            for all $x_1,x_2 \in \mathbb{R}^d$.
        \end{defn}
        Now we present the contraction mapping theorem.
        \begin{thm}{Contraction Mapping Theorem}{ContractMap}
            Let $f:\mathbb{R}^n \to \mathbb{R}^n$ be a contraction mapping. Then, the following properties hold for equation $x=f(x)$:
            \begin{itemize}
                \item Existence: There exists a fixed point $x^*$ satisfying $f(x^*)=x^*$;
                \item Uniqueness: The fixed point $x^*$ is unique;
                \item Algorithm: The following sequence $\{x_k\}$ converges to the fixed point $x^*$ exponentially, for any initial $x_0$.
                    \begin{equation}
                        x_{k+1} = f(x_k)
                    \end{equation}
            \end{itemize}
        \end{thm}
        \begin{proof}
            \textit{Part 1.} (Convergence) The sequence $\{x_k\}$ defined as $x_{k+1}=f(x_k)$ converges.

            We show that the sequence $\{x_k\}$ is a Cauchy sequence; i.e., there always exists an $N \in \mathcal{N}$ such that $\|x_m-x_n\| < \epsilon$ for any $\epsilon > 0$.
            The Cauchy sequence has a nice property; the Cauchy sequence always converges. This can be proved by using the properties of $\mathcal{R}$, that the $\mathcal{R}$ is a complete set.
            We won't show this here. The following holds:
            \begin{align}
                \|x_{k+1}-x_k\| = &\|f(x_k) - f(x_{k-1})\| \\
                                  &\leq \gamma \|x_k - x_{k-1}\| \\
                                  &\leq \gamma^2 \|x_{k-1}-x_{k-2}\| \\
                                  &\leq \gamma^3 \|x_{k-2}-x_{k-3}\| \\
                                  &\vdots \\
                                  &\leq \gamma^k \|x_1 - x_0\|
            \end{align}
            Using this, we can show that
            \begin{align}
                \|x_m - x_n\| &\leq \|x_m - x_{m-1}\| + \|x_{m-1} - x_{m-2}\| + \cdots + \|x_{n+1} - x_n\| \\
                              &\leq \gamma^{m-1}\|x_1-x_0\| + \gamma^{m-2}\|x_1-x_0\| + \dots + \gamma^{n}\|x_1-x_0\| \\
                              &= \|x_1-x_0\| (\gamma^{m-1} + \gamma^{m-2} + \dots + \gamma^n) \\
                              &= \|x_1-x_0\| \frac{\gamma^n(1-\gamma^{m-n})}{1-\gamma} \\
                              &\leq \|x_1-x_0\| \frac{\gamma^n}{1-\gamma}
            \end{align}
            where $m \geq n$, without loss of generality. Since $\gamma \in (0,1)$, the last line converges to 0. That is, the sequence $\{x_k\}$ is a Cauchy sequence.

            \textit{Part 2.} (Existence) $\lim_{k \to \infty} x_k$ is a solution of fixed-point equation $x=f(x)$.

            Let $x^* = \lim_{k \to \infty} x_k$. Then, following holds.
            \begin{align}
                \|x^*-f(x^*)\| &\leq \|x^*-x_k\| + \|f(x_k)-f(x^*)\| + \|x_k-f(x_k)\| \\
                               &\leq \|x^*-x_k\| + \gamma \|x_k - x^*\| + \|x_k - f(x_k)\|
            \end{align}
            Using the results of \textit{Part 1} and squeeze theorem, we can show that $x^*=f(x^*)$.

            \textit{Part 3.} (Uniqueness) The fixed-point $x^*$ is unique.

            Suppose there is another fixed point $x'$. Then, $\|x^*-x'\| = \|f(x^*)-f(x')\| \leq \gamma \|x^*-x'\|$. Since $\gamma \in (0,1)$, the only possibility is $x^*=x'$.


        \end{proof}

    \subsection{Existence and Uniqueness of the Solution of Bellman Equations}\label{sec:SolveBellmanOptEq}
        Using the contraction mapping theorem, we can prove the following theorem.

        \begin{thm}{Existence and Uniqueness of the Solution of Bellman Equations}{exis-uniq-gen-bellman}
            Let $T:\mathbb{R}^\mathcal{S} \to \mathbb{R}^\mathcal{S}$ be the generalized Bellman operator: $[Tv](s) = \bigotimes_a \Big( r(s,a) + \gamma \bigoplus_{s'}v(s') \Big)$.
            Then, if both $\bigotimes$ and $\bigoplus$ is non-expanding, i.e.,
            $\left|\bigotimes_a f_1(s,a) - \bigotimes_a f_2(s,a)\right| \leq \max_a |f_1(s,a) - f_2(s,a)|$, and $0 \leq \gamma < 1$, the associated Bellman equation
            \begin{equation}
                v(s) = \bigotimes_a \Big(r(s, a) + \gamma \bigoplus_{s'} v(s') \Big) 
            \end{equation} 
            always has solution and the solution is unique.
        \end{thm}
        \begin{proof}
            We show that the operator $T$ is a contraction operator, and use Theorem~\ref{thm:ContractMap}.
            Let $h_v(s,a) = r(s,a) + \gamma \bigoplus_{s'} v(s')$. Then,
            \begin{equation}
                h_{v_1}(s,a) - h_{v_2}(s,a) = \gamma \bigoplus_{s'} \left( v_1(s') - v_2(s') \right)
            \end{equation}
            and using the non-expansion assumption of $\bigoplus$,
            \begin{equation}
                |h_{v_1}(s,a) - h_{v_2}(s,a)| \leq \gamma \max_{s'}|v_1(s') - v_2(s')| = \gamma \|v_1 - v_2\|_\infty
            \end{equation}
            for all $(s, a)$.
            Now, using this and non-expansion assumption of $\bigotimes$, we can get
            \begin{equation}
                |[Tv_1](s) - [Tv_2](s)| = |\bigotimes_a h_{v_1}(s,a) - \bigotimes_a h_{v_2}(s, a)| \leq \max_a |h_{v_1}(s,a) - h_{v_2}(s,a)| \leq \gamma \|v_1 - v_2\|_\infty
            \end{equation}
            for all $s$. Therefore,
            \begin{equation}
                \|[Tv_1](s) - [Tv_2](s)\|_\infty \leq \gamma \|v_1 - v_2\|_\infty
            \end{equation}.
        \end{proof}

        Now, if we prove that $max_a$, $\sum_{s'} \pi(s,a)$ and $\sum_{s' \in \mathcal{S}} p(s'|s,a) v(s')$ is a non-expansion operator,
        the existence and uniqueness of the solution of Bellman equation and Bellman optimality equation is guaranteed, repectively.
        TODO: Show it!

\section{Optimal Policy, Optimal State Value and Bellman Optimality Equation}\label{sec:OptPolOptStateValBOE}
    Now, we should prove that the solution of the Bellman optimality equation is really an optimal state-value function,
    and prove that the greedy policy derived from it is the optimal policy.
    \begin{thm}{Optimality of $v^*$ and $\pi^*$}{OptPolOptStateVal}
        Let $v^*$ be the solution of BOE, and let $\pi^* = \arg\max_{\pi \in \Pi}(r_\pi + \gamma P_\pi v^*)$.
        The, the $v^*$ is an optimal state vaule function, and the $\pi^*$ is an optimal policy. That is, for any policy $\pi$,
        following holds;
        \begin{equation}
            v^* = v_{\pi^*} \geq v_\pi
        \end{equation}
        where $v_\pi$ is the state value function of $\pi$, and $\geq$ is applied elementwise.
    \end{thm}
    \begin{proof}
        For any policy $\pi$, $v_\pi = r_\pi + \gamma P_\pi v_\pi$ holds (Bellman equation). Then,
        \begin{align}
            v^* &= \max_{\tau \in \Pi}(r_\tau + \gamma P_\tau v^*) \\
                &= r_{\pi^*} + \gamma P_{\pi^*} v^* \\
                &\geq r_\pi + \gamma P_\pi v^*
        \end{align}
        and we have
        \begin{equation}
            v^*-v_\pi \geq \gamma P_\pi(v^*-v_\pi).
        \end{equation}
        Using this inequality repeatedly, we have $v^*-v_\pi \geq \gamma^n P_\pi^n (v^*-v_\pi)$ for any $n \in \mathbb{N}$.
        Since $\gamma \in (0,1)$ and $P_\pi$ is a nonnegative matrix with all its elements less than or equal to 1,
        $\lim_{n \to \infty} \gamma^n P_\pi^n (v^*-v_\pi)=0$. That is, $v^*-v_\pi \geq \lim_{n \to \infty} \gamma^n P_\pi^n (v^*-v_\pi)=0$.
    \end{proof}
    The following theorem examines the optimal policy more precisely.
    \begin{thm}{Greedy Optimal Policy}{GreedyOptPol}
        For any $s \in \mathcal{S}$, the following deterministic policy is an optimal policy for solving the BOE.
        \begin{equation}
            \pi^*(a|s) = 
            \begin{cases}
                1, &a = \arg\max_a q^*(a,s) \\
                0, &a \neq \arg\max_a q^*(a,s)
            \end{cases}
        \end{equation}
        where $q^*(s,a) = \sum_{r \in \mathcal{R}}p(r|s,a)r + \gamma \sum_{s' \in \mathcal{S}}p(s'|s,a)v^*(s')$.
    \end{thm}
    \begin{proof}
        From Theorem~\ref{thm:SolutionBellmanOptEq} and~\ref{defn:BellOptEq}, we see that
        \begin{align}
            v^*(s) &= \max_{\pi \in \Pi}\sum_{a \in \mathcal{A}}\pi(a|s)q^*(s,a) \\
                        &\leq \left(\max_{\pi \in \Pi}\sum_{a \in \mathcal{A}}\pi(a|s)\right) \max_{a'} q^*(s,a') \label{tmpeq:2} \\
                        &= \max_{a'}q^*(s,a')
        \end{align}
        and the given deterministic policy $\pi^*$ satisfies the equality of inequality~\ref{tmpeq:2}.
    \end{proof}
    The important truth is that the optimal state value is unique but the optimal policy is not. The Theorem~\ref{thm:GreedyOptPol}
    only shows that there always exists at least one deterministic optimal policy. There still can be other optimal policies, too, stochastic or whatever.

\section{Factors that Influence Optimal Policies}
    Now, we analyze the factors of BOE that influences optimal policies. The factors are: reward and discount rate.
    \subsection{Reward}
        Changing reward can lead to slightly different behavior of policies. For example, in the Figure 3.4(a) of \cite{Zhao2025} (gridworld), one can see that
        the agent tries to move to the goal even though there are forbidden areas. The forbidden area gives the negative reward, but the cumulative reward is greater than
        going around the forbidden areas. This is why reward (signal) is important.
    \subsection{Discount Rate}
        High dicount rate leads to far-sighted agent, and low discount rate leads to short-sighted agent. In the extreme, when $\gamma$, the agent only
        take into account the immediate reward.
    \subsection{Invariance of the Optimal Policy}
        Changing the reward may lead to different optimal policies, but this is not always the case. If we apply the same affine transformation
        to all rewards, the optimal policy remains the same.
        \begin{thm}{Optimal Policy Invariance}{OptPolInvariance}
            Let $v^* \in \mathbb{R}^{|\mathcal{S}|}$ be the optimal state value satisfying $v^* = \max_{\pi \in \Pi}{(r_\pi + \gamma P_\pi v^*)}$.
            If every reward $r \in \mathcal{R}$ goes through the same affine transformation $\alpha * r + \beta$, where $\alpha, \beta \in \mathbb{R}$ and $\alpha > 0$,
            the corresponding optimal state value $v'$ is also an affine transformation of $v^*$:
            \begin{equation}
                v' = \alpha v^* + \frac{\beta}{1 - \gamma}\mathbb{1}
            \end{equation}
            where $\gamma \in (0,1)$ is the discount rate and $\mathbb{1} = [1,\dots,1]^\top$.
            Due to this property, the optimal policy stays invariant.
        \end{thm}
    \chapter{Stochastic Approximation Theory}\label{ch:stochasticApprox}
    \section{Introduction}
This chapter introduces measure-theoretic probability, convergence and several stochastic approximation theorems, which is a theoretic foundation for TD-Learning algorithms.
    \subsection{Prerequisite}
        The subsections ahead requries a basic knowledge about measure-theoretic probability, which are provided in \ref{ch:Prob}.
        The probability theory was re-established by the measure theory, which enabled to define all the ambiguous concepts of probability theory rigorously.
        For comprehensive introduction about measure theory, see \cite{Axler2020} and for more information about measure-theoretic probability, see \cite{Brzezniak2007}.
    \subsection{Robbins-Monro Algorithm}
        Suppose we want to find the root of the equation
        \begin{equation}
            g(w) = 0.
        \end{equation}
        There are many other iterative algorithms for finding the root of a given equation. Then, why do we need
        stochastic approximation? In real world problems, the expression of function $g$ or its derivative is unknown,
        and we may even only be able to obtained the noisy observations of $g$. That is, with the original $g$ and noise $\eta$,
        we can only observe $\tilde{g}(w,\eta) = g(w) + \eta$. The SGD method is one of the most well-known solution to problems like this.
        But here, we introduce a more general algorithms of stochastic approximation, called \textbf{Robbins-Monro algorithm}.

        \begin{thm}{Robbins-Monro Algorithm}{RMalg}
            Let $g$ be a function and let $\tilde{g}(w, \eta) = g(w) + \eta$ where $\eta$ is a random noise.
            Then, the sequence $\{w_k\}$ defined as $w_{k+1} = w_k - a_k \tilde{g}(w_k, \eta_k)$
            where $\{a_k\}$ is a sequence of positive coefficients almost surely converges to the root of equation $g(w)=0$
            under following conditions:
            \begin{enumerate}
                \item $0 < c_1 \leq \nabla_w g(w) \leq c_2$ for all $w$;
                \item $\sum_{k=1}^{\infty}a_k = \infty$ and $\sum_{k=1}^{\infty}a_k^2 < \infty$;
                \item $\E[\eta_k|\mathcal{H}_k]=0$ and $\E[\eta_k^2|\mathcal{H}_k] < \infty$; 
            \end{enumerate}
            where $\mathcal{H}_k = \{w_k, w_{k-1}, \dots\}$.
        \end{thm}
        \begin{proof}
        TODO: proof of RM algorithm
        \end{proof}
\part{Algorithms}
    \chapter{Dynamic Programming}
    \section{Introduction}
This chapter discuss and dyamic programming algorithms for finding optimal policies.
Since these are dynamic prgramming algorithms, a perfect system model is required; this mean that
it is rarely used in practice, but these algorithms important foundations of the model-free algorithms.
Three algorithms will be presented;
\begin{itemize}
    \item Value Iteration: The algorithms suggested by the contraction mapping theorem \ref{thm:SolutionBellmanOptEq};
    \item Policy Iteration:
    \item Trucated Policy Iteration: The unified algorithm which includes value iteration and policy iteration as its special cases.
\end{itemize}

\section{Value Iteration}
    \subsection{Main Idea}
        The main idea of \textbf{value iteration} algorithms is the algorithms suggested in Theorem \ref{thm:SolutionBellmanOptEq}.
        In particular, the algorithms is
        \begin{equation}
            v_{k+1} = \max_{\pi \in \Pi}{(r_\pi + \gamma P_\pi v_k)},\quad k=0,1,2,\dots
        \end{equation}
        Breaking down, this iterative algorithm has two steps;
        \begin{enumerate}
            \item \emph{Policy Update}: The algorithm finds a policy $\pi_{k+1}$ that solves following optimization problem;
                \begin{equation}
                    \arg\max_\pi(r_\pi + \gamma P_\pi v_k)
                \end{equation}
            \item \emph{Value Update}: With the policy $\pi_{k+1}$, the algorithm calculates a new state value function $v_{k+1}$;
                \begin{equation}
                    v_{k+1} = r_{\pi_{k+1}} + \gamma P_{\pi_{k+1}} v_k
                \end{equation} and this $v_{k+1}$ is used in next iteration.
        \end{enumerate}
    \subsection{Pseudocode}
        \Algorithm{valueIter}
        {Value Iteration Algorithm}
        {\inputitem{$p(s',r|s,a)$}{the dynamics of the system} \\
         \inputitem{$v$}{initial guess}}
        {\inputitem{$v^*,\pi^*$}{optimal state value and optimal policy}}
        {
            $\pi \leftarrow$ greedy policy derived from $v$\;
            \Repeat{convergence}
            {
                \For{$s \in \mathcal{S}$} {
                    \For{$a \in \mathcal{A}(s)$} {
                        $q(s,a) \leftarrow \sum_r p(r|s,a)r + \gamma \sum_{s'}p(s'|s,a)v(s')$\;
                    }
                    $a^* \leftarrow \arg\max_a q(s,a)$\;
                    \tcp{policy update}
                    $\pi(s,a)=
                    \begin{cases}
                        1 &\text{if $a$ = $a^*$} \\
                        0 &\text{otherwise}
                    \end{cases}$\;
                    \tcp{value update}
                    $v(s) = \max_a{q(s,a)}$\;
                }
            }
            \Return{$v$, $\pi$}
        }
        Note that the value update of Algorithm \ref{alg:valueIter} is simply substituting the original state values with
        maximum action values since we are using greedy policy.
    \subsection{Convergence}
        The convergence can be proved with Theorem \ref{thm:ContractMap} and Theorem \ref{thm:ContPropBellmanOptEq}.
\section{Policy Iteration}
    \subsection{Main Idea}
        The \textbf{policy iteration} is based on the policy improvement theorem. 
        Like value iteration, the policy iteration consists of two steps:
        \begin{enumerate}
            \item \emph{Policy Evaluation}: For a given policy $\pi_k$, we calculate the corresponding state values by solving following Bellman equation:
                \begin{equation}
                    v_{\pi_k} = r_{\pi_k} + \gamma P_{\pi_k}v_{\pi_k}
                \end{equation}
            \item \emph{Policy Improvement}: Using the computed state value $v_k$, we obtain the improved policy as
                \begin{equation}
                    \pi_{k+1} = \arg\max_\pi(r_\pi + \gamma P_\pi v_{\pi_k})
                \end{equation} and here is where the policy improvement theorem is applied.
        \end{enumerate}
        The difference is that the value iteration is that instead of taking a state value and computing the policy from it,
        the policy iteration takes a policy and compute the value corresponding to the policy.

        Also, note that policy iteration has another iterative algorithm inside it: the policy evaluation uses iterative algorithm for solving the
        Bellman equation $v_{\pi_k} = r_{\pi_k} + \gamma P_{\pi_k}v_{\pi_k}$.

        
    \chapter{Monte Carlo Methods}
    \section{Introduction}
This chapter introduces a \textbf{Monte Carlo methods} for reinforcement learning.
These algorithms all starts from trying to obtain the optimal policy by solving the Bellman optimality equation.
The difference of pure dynamic programming algorithms and Monte Carlo methods is that the latter
do not require a perfect dynamics of the model, i.e., it is \emph{model-free}.

Then, how do we get the optimal policy without a perfect model? We use the data that the model gives us.
Using the data, we can convert the policy iteration algorithm~\ref{alg:policyIter} to the model-free version.
Specifically, we estimate the action value $q(s,a)=\E[G_t|S_t=a,A_t=a]$ using Monte Carlo sampling.
\section{Monte Carlo Basic}
    \subsection{Main Idea}
        The \textbf{Monte Carlo basic} is the simplest Monte Carlo method algorithm.
        Suppose that we run the agent starting from state-action pair $(s,a)$, following policy $\pi$, and let $g_\pi^{(i)}(s,a)$
        be the return of the episode. Then, we can approximate the action value $q_\pi(s,a)$ as following:
        \begin{equation}
            q_\pi(s,a) = \E[G_t|S_t=s,A_t=a] \approx \frac{1}{n} \sum_{i=1}^n g_\pi^{(i)}(s,a).
        \end{equation}
        If the number $n$ is sufficiently large, the approximation will be sufficiently accurate. This can be guaranteed due to the law of large numbers.
    \subsection{Pseudocode}
        \Algorithm{mcBasic}
        {Monte Carlo Basic}
        {
            \inputitem{$\pi_0$}{initial guess}
        }
        {
            \inputitem{$\pi^*$}{optimal policy}
        }
        {
            \Repeat{convergence of $\pi$}
            {
                \ForEach{$s \in \mathcal{S}$}
                {
                    \ForEach{$a \in \mathcal{A}(s)$}
                    {
                        \tcp{policy evaluation}
                        collect sufficiently many episodes starting from $(s,a)$, following policy $\pi_k$\;
                        $q_{\pi_k}(s,a) \leftarrow$ the average return of all episodes starting from $(s,a)$\;
                    }
                }
                \tcp{policy improvement}
                $\pi_{k+1}(a|s) = \begin{cases}
                    1 &\text{if $a=\arg\max_{a'}{q_{\pi_k}(s,a')}$} \\
                    0 &\text{otherwise}
                \end{cases}$\;
            }
            \Return{$\pi$}
        }
    \subsection{Convergence}
    \subsection{Pros and Cons}
        Cons:
        \begin{itemize}
            \item Very inefficient sample usage. Due to this property, MC-Basic is not used in practical.
                To deal with this, we use MC with exploring start, which is presented in next section.
        \end{itemize}
\section{Monte Carlo with Exploring Starts}
    \subsection{Main Idea}
        The MC-Basic algorithms is not sample-efficient, since it only updates only q value of
        one state-action pair each episode. To deal with this low efficiency, we use all the
        \emph{visit} to the state-action pair in each episode.
        
        Suppose that the agent has gone through following trajectory:
        \begin{equation}
            s_1 \xrightarrow{a_1} s_2 \xrightarrow{a_2} s_3 \xrightarrow{a_3} s_4 \xrightarrow{a_4} \dots
        \end{equation}
        To increase the sample efficiency, we use the visits to the state-action pairs in the trajectory.
        That is, using the above episode, we not only use $(s_1,a_1)$ for estimating $q(s_1,a_1)$ but use $(s_2,a_2)$ for estimating $q(s_2,a_2)$, and so on.
        There are two different ways to do this:
        \begin{itemize}
            \item The first way is to only use the state-action pairs when it is first encountered. This is called \emph{first-visit method};
            \item The second way is to use all the visit to the state-action pairs. This is called \emph{every-visit method}.
        \end{itemize}

        Also, we update the estimate of q values after each episode ends. This falls into the scope of \emph{generalized policy iteration},
        since we are using the value approximation even it is not so accurate.
    \subsection{Pseudocode}
        \Algorithm{mc-explore}
        {Monte Carlo with Exploring Starts (every-visit)}
        {
            \inputitem{$\pi_0$}{initial policy}
        }
        {
            \inputitem{$\pi^*$}{optimal policy}
        }
        {
            \tcp{initialize}
            \ForEach{valid state-action pair $(s,a)$}
            {
                $\pi \leftarrow \pi_0$\;
                $q(s,a) \leftarrow$ initial value\;
                \tcp{holds the sum of returns of (sub)episode starting from $(s,a)$}
                $\text{Ret}(s,a) \leftarrow 0$\;
                \tcp{counts the number of visits of $(s,a)$}
                $\text{N}(s,a) \leftarrow 0$\;
            }
            \Repeat{convergence}
            {
                select a starting valid state-action pair $(s,a)$, ensuring that all the pairs can be possibly selected\;
                generate an episode starting from $(s,a)$, following current policy $\pi$: $s_0, a_0, r_1, s_1, a_1, \dots, s_{T-1}, a_{T-1}, r_T$\;
                $G \leftarrow 0$\;
                \For{$t=T-1,T-2,\dots,0$}
                {
                    $G \leftarrow r_{t+1} + \gamma G$\;
                    $\text{Ret}(s,a) \leftarrow \text{Ret}(s,a) + G$\;
                    $\text{N}(s,a) \leftarrow \text{N}(s,a) + 1$\;
                    \tcp{policy evaluation}
                    $q(s_t,a_t) \leftarrow \frac{\text{Ret}(s,a)}{\text{N}(s,a)}$\;
                    \tcp{policy improvement}
                    $\pi_{k+1}(a|s) = \begin{cases}
                        1 &\text{if $a=\arg\max_{a'}{q_{\pi_k}(s,a')}$} \\
                        0 &\text{otherwise}
                    \end{cases}$\;
                }
            }
        }
        
    \subsection{Implementation Details}
        The exploring-start condition can be achieved by random start state.

    \subsection{Pros and Cons}
    Pros:
    \begin{itemize}
        \item Better sample efficiency then MC-Basic;
        \item All state-action pairs can be visited due to the exploring-start, and this makes the estimation more accurate.
    \end{itemize}
    Cons:
    \begin{itemize}
        \item Only applicable to episodic tasks, which is a common drawback for all MC methods. This property leads to following problem, too;
        \item High variance since the return holds all the noises and randomnesses, and the return is used for update;
        \item If the exploring start condition cannot be satisfied, the MC with Exploring Starts cannot be used. This is common in real-world problems.
        To deal with this, another algorithms without exploring starts will be introduced next section;
    \end{itemize}
    \chapter{Temporal Difference Methods}
    \section{Introduction}
    This chapter presents the temporal-difference (TD) methods for reinforcement learning.
    The main difference of MC and TD is that TD learning is incremental, while MC is not.
    That is, TD methods are appliable on continuous tasks, unlike MC method.
\section{TD Learning of State Values}
    The general algorithm of TD method with state value is given in this section.
    In practice, we use TD method for action values, but the one given here is a foundation of
    them.
% Fill here later.
\section{Sarsa}
    Now, we introduce the most basic TD-algorithm, the \textbf{Sarsa}. Its name originates from the
    tuple $(s,a,r,s',a')$ it uses.
    
    In the last section, the TD algorithm of learning state value was mainly discussed. Here, we directly learn the
    action values instead of state values, since it can be combined with a policy improvement step to learn optimal policies.
    \subsection{Main Idea}
        Suppose we have trajectory generated from policy $\pi$: $(s_0,a_0,r_1,s_1,a_1,\dots,s_t,a_t,r_{t+1},s+{t+1},a_{t+1},\dots).$
        Note that the trajectory does not have to end. We use following TD-update every step for estimating the action values:
        \begin{align}
            q_{t+1}(s_t,a_t) &= q_t(s_t,a_t) + \alpha_t(s_t,a_t)\Big[ r_{t+1} + \gamma q_t(s_{t+1},a_{t+1}) - q_t(s_t,a_t) \Big] \label{eq:SarsaUpdate}\\
            q_{t+1}(s,a) &= q_t(s,a),\text{ for all }(s,a) \neq (s_t,a_t)
        \end{align}    

        In fact, the Sarsa solves following Bellman equation in a stochastic approximation manner.
        \begin{equation}
            q_\pi(s,a) = \E[R + \gamma q_\pi(S',A')|s,a] \text{ for all }(s, a) \label{eq:sarsaBellman}
        \end{equation}
        Equation \ref{eq:sarsaBellman} is the Bellman equation about action values. The proof is given below.
        \begin{proof}
            Note that
            \begin{align}
                p(s',a'|s,a) &= p(s'|s,a)p(a'|s',s,a) \\
                             &= p(s'|s,a)p(a'|s') \text{ (due to the Markov property)} \\
                             &= p(s'|s,a)\pi(a'|s').
            \end{align}
            Using this, the equation of Theorem \ref{thm:BellmanEq-AV} can be written as
            \begin{equation}
                q_\pi(s,a) = \sum_{r}rp(r|s,a) + \gamma \sum_{s',a'}q_\pi(s',a')p(s',a'|s,a)
            \end{equation}
        \end{proof}
    \subsection{Pseudocode}
        \Algorithm{sarsa}
        {Sarsa}
        {\inputitem{$\alpha_t(s,a) \in (0,1]$}{step size parameter satisfying RM condition}}
        {\inputitem{$q_*$}{the optimal policy}}
        {
            \tcp{initialize}
            \ForEach{$(s,a) \in \mathcal{S} \times \mathcal{A}$}
            {$q(s,a) \leftarrow
            \begin{cases}
                \text{arbitrary,} &\text{ if $s \in \mathcal{S}^+$ and $a \in \mathcal{A}(s)$} \\
                0 &\text{ if $s \in \mathcal{S} \setminus \mathcal{S}^+$}
            \end{cases}$\;}
            $\pi \leftarrow$ policy derived from $q$ (e.g. $\varepsilon$-greedy, softmax, GLIE)\;

            \ForEach{episodes}
            {
                initialize $s$\;
                choose action $a$ from $s$ using policy $\pi$\;
                \Repeat{s is in terminal state}
                {
                    take action $a$, observe $r$, $s'$\;
                    choose action $a'$ from $s'$ using policy derived from $\pi$\;
                    \tcp{update q value}
                    $q(s,a) \leftarrow q(s,a) + \alpha_t(s,a) \big[ r + \gamma q(s',a') - q(s,a) \big]$\;
                    \tcp{update policy}
                    update $\pi$ according to $q(s,a)$
                    $s \leftarrow s'$\;
                    $a \leftarrow a'$\;
                }
            }
            \Return{$q \simeq q_*$}
        }
    \subsection{Implementation Detail}
        I usually intialize q value without minding if the state is terminal or not. While traning the agent, when the agent meets terminal state,
        I just check it in place and set the $q_t(s_{t+1},a_{t+1}) = 0$. Actually I don't even change the q value of terminal state to 0. I just leave it untouched.
    \subsection{Convergence}
        \begin{thm}{Convergence of Sarsa}{convSarsa}
            Given a policy $\pi$, by the Sarsa algorithm \ref{eq:SarsaUpdate}, $q_t(s,a)$ converges almost surely to the action value $q_\pi(s,a)$ as $t \to \infty$ for all
            $(s,a)$ if $\sum_t \alpha_t(s,a) = \infty$ and $\sum_t \alpha_t^2(s,a) < \infty$ for all $(s,a)$.
        \end{thm}
        \begin{proof}
            TODO: organize the proof!
        \end{proof}
    \subsection{Pros and Cons}
        Pros:
        \begin{itemize}
            \item It learns safer paths by accounting for penalties from exploratory moves. This can be checked with cliff walking of \cite{Sutton2020}, with comparison to Q-Learning.
        \end{itemize}
        Cons:
        \begin{itemize}
            \item It learns suboptimal policies first, due to the behavior of searching for safe paths in the initial training process.
            \item The variance can be high, since in the initial training process the q values are very noisy and inaccurate, and the agent use the q value of it directly in the update.
                This can be dealt with the Expected Sarsa \ref{sec:ExpectedSarsa}.
        \end{itemize}
\section{Q-Learning}
    In this section we learn another TD-learning algorithm, the \textbf{Q-Learning}.
    Unlike Sarsa, Q-Learning bootstraps the q value using the maximum action value of current state.
    \subsection{Main Idea}
        The Q-Learning uses following update equation for action values:
        \begin{align}
            q_{t+1}(s_t,a_t) &= q_t(s_t,a_t) + \alpha_t(s_t,a_t)\Big[ r_{t+1} + \gamma \max_{a \in \mathcal{A}(s_{t+1},a)}{q_t(s_{t+1},a)} - q_t(s_t,a_t) \Big] \\
            q_{t+1}(s,a) &= q_t(s,a),\text{ for all }(s,a) \neq (s_t,a_t)
        \end{align}

        The Q-Learning solves following Bellman optimality equation:
        \begin{equation}
            q(s,a) = \E\Big[ R_{t+1} + \gamma \max_aq(S_{t+1},a)|S_t=s,A_t=a \Big]
        \end{equation}
        \begin{proof}
            Check for \cite{Zhao2025} pg.140.
        \end{proof}

        \begin{rem}{On-policy and Off-policy}{onPoloffPol}
            As one can see, the Sarsa solves Bellman equation and then make policy update, but Q-Learning directly learns the optimal policy.
            The \textbf{behavior policy}, is used to make the trajectory. The \textbf{target policy} is the policy that learns the optimal policy.
            When the behavior policy and target policy equals, the algorithms is called \textbf{on-policy} algorithm. If not, it is called \textbf{off-policy}.
            Off policy can separate the exploration with target policy. The target policy only has to care about choosing the best action in current situation.
            Since the off-policy uses policy other than its target policy, it requires some techniques such as \textbf{importance sampling} and \textbf{tree-backup} algorithms.
            These will be covered later.
            
            The on-policy algorithms cannot be implemented as an \textbf{offline} fashion, since it requires real-time interaction with environment.
            The off-policy algorithms, in contrast, can be implemented as an offline or online fashion.
        \end{rem}

    \subsection{Pseudocode}
        \Algorithm{qlearning}
        {Q-Learning}
        {
            \inputitem{$\alpha_t(s,a) \in (0,1]$}{step size parameter satisfying RM condition} \\
            \inputitem{$b$}{behavior policy} \\
            \inputitem{$\pi$}{target policy}
        }
        {
            \inputitem{$q_*$}{the optimal policy}
        }
        {
            \tcp{initialize}
            \ForEach{$(s,a) \in \mathcal{S} \times \mathcal{A}$}
            {$q(s,a) \leftarrow
            \begin{cases}
                \text{arbitrary,} &\text{ if $s$ is not in terminal states} \\
                0 &\text{ if $s$ is in terminal states}
            \end{cases}$\;}
            $s \leftarrow$ initialize\;
            \ForEach{episode} {
                \Repeat{$s$ is not terminal state} {
                    choose action $a$ from $s$ using behavior policy $b$\;
                    observe state $s'$ and reward $r$\;
                    \tcp{update q value}
                    $q(s,a) \leftarrow q(s,a) + \alpha_t(s,a)\Big[ r + \gamma \max_a{q(s',a)} - q(s,a) \Big]$
                    \tcp{update target policy}
                    update target policy according to $q$\;
                    $s \leftarrow s'$\;
                }
            }
        }
    \subsection{Implementation Detail}
        The behavior policy can be any policy that explores well, but in practice we use decaying $\varepsilon$-greedy due to the convergence rate. (I guess...)
    \subsection{Convergence}
        TODO: Convergence of Q-Learning
    \subsection{Pros and Cons}
        Pros:
        \begin{itemize}
            \item It is off-policy. Whatever the behavior policy does, the target policy learns the optimal policy.
            \item Does not fears getting penalty, as opposed to Sarsa. This can be checked in \cite{Sutton2020}, cliff walking.
        \end{itemize}
        Cons:
        \begin{itemize}
            \item Can suffer from \textbf{maximization bias}, since the Q-Learning uses maximum q value for bootstrapping. That is, the agent keeps
                overestimate the q values, which can lead to slow convergence. For more information, see \cite{Sutton2020}. To cope with this, we can use \textbf{Double Q-Learning}.
        \end{itemize}
\section{Expected Sarsa} \label{sec:ExpectedSarsa}
    \chapter{Function Approximation Methods}
    \section{Introduction}
In this chapter, value function approximation methods will be covered. Unlike the tabular methods, the approximate method tries to
approximate the state value or action value of state-action pairs, using a function. This solves the curse of dimensionality, which is a critical problem for tabular methods.

But there are tradeoffs, too. For example, the convergence property breaks down easily. When off-policy, function approximation, and bootstrapping meets together,
(\textbf{the deadly triad}, \cite{Sutton2020}), the algorithm lost its convergence property and fails to find optimal policy. There are various ways to solve this problem.
Some solve this with practical methods (replay buffer, PER, double learning, etc.), and some solve this theoretically (GTD2, Emphatic-TD, etc.).
I'll try to cover both of them in this note.

\section{TD Learning of State Values based on Function Approximation}
Methods of estimating the state value with a function approximation is introduced in this section.
    \subsection{Objective Function}
        Let $v_\pi(s)$ be the true state value function and let $\hat{v}_\pi(s, w)$ be the approximation of it with a parameter $w$.
        To find the optimal parameter $w$ which best approximates the true state value function, we minimize the objective function. 
        In particuluar, the objective function is given as
        \begin{equation}
            J(w) = \E[(v_\pi(S) - \hat{v}(S, w))^2],
        \end{equation}
        where $S$ is a random variable of state space. Then what is the distribution of a random variable $S$?
        We can use \emph{stationary distribution} of Markov Process. Stationary distribution of a Markov Process describes the long-term behavior of a Markov decision process.
        That is, the probability of agent being located at certain state after a long execution can be described by this distribution.

        Let $d_\pi(s)$ be the stationary distribution of the Markov process under policy $\pi$. Then, the objective functon can be written as
        \begin{equation}
            J(w) = \sum_{s \in \mathcal{S}}d_\pi(s)(v_\pi(s) - \hat{v}(s, w))^2 \label{eq:stationaryDist}
        \end{equation}
        TODO: Why stationary distribution? + Tsitsiklis and van Roy + Chapter 6

\section{Deep Q-Learning}
    \subsection{Main Idea}
        The \textbf{Deep Q-Learning} \cite{Mnih2013} is one of the earliest and successful deep reinforcement learning algorithms.
        Suprisingly, DQN algorithm are not guaranteed to converge, since the deadly triad perfectly fits into the algorithms' condition:
        it bootstraps, it is off-policy, and it uses function approximation. But there are practical methods for calming down this unstability.
        But still note that the algorithm is not guaranteed to converge.

    \subsection{Objective Function}
        The objective function which DQN minimizes is given as following:
        \begin{equation}
            J = \E\Big[\Big( R + \gamma \max_{a \in \mathcal{A}(S')}{\hat{q}(S',a,w)} - \hat{q}(S, A, w) \Big)^2\Big],
        \end{equation}
        where $\hat{q}$ is the approximation function of action values
        and $(S, A, R, S')$ are random variables denoting state, action, reward and next state, resepectively.
        Note that the squared term is a Bellman optimality error.

        Now, we can caculate the gradient and minimize the loss with gradient descents... but take a look. Unfortunately, the
        max operation is not differentiable. To resolve this, we freeze the parameter of a network inside the max operation for a short time. Specifically,
        we introduce two networks: \emph{target network and main network}. Using the main network $\hat{q}(s, a, w)$ and $\hat{q}(s, a, w_T)$,
        we can re-write the objective function as following:
        \begin{equation}
            J = \E\Big[\Big( R + \gamma \max_{a \in \mathcal{A}(S')}{\hat{q}(S',a,w_T)} - \hat{q}(S, A, w) \Big)^2\Big], \label{eq:DQNObjective}
        \end{equation}
        and the gradient can be computed about $w$. The parameters of target network and main network are synchronized regularly.

        There is one more technique required for DQN to be trained. The \emph{experience replay} collects the $(s, a, r, s')$ pair, and then
        samples the experiences with uniform random distrubution. Why is this technique required?
        
        \begin{itemize}
            \item \cite{Zhao2025} says: Compare equation \ref{eq:DQNObjective} and \ref{eq:stationaryDist}. Since we don't know the probability distribution of $(S, A, R, S')$,
                we just simply set it as uniform distrubution. However, the state-action pairs are are strongly correlated, and not uniformly distributed in single trajectory.
                To break the correlation, we uniformly draw samples from the replay buffer.
        \end{itemize}
        
        Also, since we can reuse the experiences, the sample efficiency naturally increases.
    \subsection{Pseudocode}
        \Algorithm{DQN}
        {Deep Q-Learning}
        {
            \inputitem{$b$}{behavior policy} \\
            \inputitem{$C$}{sync rate}
        }
        {\inputitem{$\hat{q}(s, a, w)$}{approximated action values}}
        {
            store the samples generated by $b$ in a replay buffer $\mathcal{B} = \{(s, a, r, s')\}$\;
            \ForEach{iteration}
            {
                uniformly draw a mini batch $B$ of samples from $\mathcal{B}$\;
                $L \leftarrow 0$\;
                \ForEach{sample $(s, a, r, s') \in B$}
                {
                    $y_T \leftarrow r + \gamma \max_{a \in \mathcal{A}(s')}{\hat{q}(s', a, w_T)}$\;
                    $L \leftarrow L + \Big( y_T - \hat{q}(s, a, w) \Big)^2$\;
                }
                minimize $L$ by updating parameter $w$\;
                set $w_t = w$ every $C$ iterations
            }
            \Return{$\hat{q}$}
            
        }
    \subsection{Implementation Details}
        Since the original DQN algorithm is highly unstable, there are several techniques to stabilize the algorithm.
        \begin{itemize}
            \item Priortized experience replay (PER)
            \item Double DQN
            \item Dueling DQN
            \item Distributional DQN
        \end{itemize}
    \subsection{Pros and Cons}
        Pros:
        \begin{itemize}
            \item 
        \end{itemize}
        Cons:
        \begin{itemize}
            \item Catastrophic forgetting: This is a common problems in neural networks.
            \item No convergence guaranteed: Note that DQN satisfies the deadly triad condition.
                And also, the algorithm does not even use the true gradient, since the target network is freezed for a moment. This is called a \textbf{semi-gradient} method,
                and no convergence is guaranteed for this method. For more information see \cite{Sutton2020}.
        \end{itemize}
    \chapter{Policy-based Methods}
    \section{Introduction}
    \textbf{Policy gradient method}, unlike value-based methods, tries to predict the optimal policy itself,
    instead of solving the Bellman equation. It has many advantages over value-based methods;
    It is more efficient in handling large state-action spaces, and has stronger generalization abilities.
\section{Metrics for Defining Optimal Policies}\label{sec:Metrics}
    The value-based methods uses state value and action value to define the optimal policies.
    New metrics for defining optimal polcies are required for policy gradient methods, since they does not use the state and action values.

    There are two metrics given in \cite{Zhao2025}.
    \begin{enumerate}
        \item average state value: $\bar{v}_\pi = \sum_{s\in\mathcal{S}}d(s)v_\pi(s) = \E_{S \sim d}[v_\pi(S)]$ where $d(s)$ is the weight of state $s$;
        \item average reward: $\bar{r}_\pi = \sum_{s\in\mathcal{S}}d_\pi(s)r_\pi(s) = \E_{S \sim d_\pi}[r_\pi(S)]$ where $d_\pi$ is the stationary distribution and $r_\pi(s) = \sum_{a \in \mathcal{A}}\pi(a|s,\theta)r(s,a)$.
    \end{enumerate}
    Actually the two metrics are equivalent in discounted case $\gamma < 1$. In particular, $\bar{r}_\pi = (1-\gamma)\bar{v}_\pi$. Proofs can be found at \cite{Zhao2025}.

    \subsection{Average State Value}
        Average state value has two equivalent expressions.
        \begin{itemize}
            \item Suppose the agent has collected rewards $\{R_{t+1}\}_{t=0}^\infty$. Then,
            \begin{align}
                \E[\sum_{t=0}^\infty \gamma^t R_{t+1}] &= \sum_{s \in \mathcal{S}}d(s)\E\left[\sum_{t=0}^\infty \gamma^t R_{t+1} \mid| S_0 = s\right] \\
                                                       &= \sum_{s \in \mathcal{S}}d(s)v_\pi(s) \\
                                                       &= \bar{v}_\pi
            \end{align}
            \item $\bar{v}_\pi$ can be rewritten as the inner product of two vectors, $v_\pi = [\dots,v_\pi(s),\dots]$ and $d = [\dots,d(s),\dots]$.
                \begin{equation}
                    \bar{v}_\pi = d^\top v_\pi
                \end{equation}
        \end{itemize}

    \subsection{Average Reward}
        Average reward has two equivalent expressions.
        \begin{itemize}
            \item Suppose the agent has collected rewards $\{R_{t+1}\}_{t=0}^\infty$. Then,
                \begin{equation}
                    \lim_{n \to \infty}\frac{1}{n}\E\left[ \sum_{t=0}^{n-1}R_{t+1} \right] = \bar{r}_\pi
                \end{equation}
                The proof can be found at \cite{Zhao2025}.
            \item Same as before, $\bar{r}_\pi$ can be rewritten as the inner product of two vectors.
        \end{itemize}

\section{Policy Gradient Theorem}
    The policy gradient method mainly depends on the \textbf{policy gradient theorem}.
    \begin{thm}{Policy Gradient Theorem}{policyGrad}
        Let $J(\theta)$ be one of the metrics proposed in Section~\ref{sec:Metrics}. Then, the gradiento of $J(\theta)$ is
        \begin{equation}
            \nabla_\theta J(\theta) = \sum_{s\in\mathcal{S}}\eta(s)\sum_{a\in\mathcal{A}(s)}\nabla_\theta \pi_\theta(a|s)q_{\pi_\theta}(s,a)
        \end{equation}
        where $\eta$ is a state distrubution. Moreover, the equivalent equation expressed in terms of expectation exists:
        \begin{equation}
            \nabla_\theta J(\theta) = \E_{S \sim \eta, A \sim \pi_\theta(\cdot|S)} \Big[ \nabla_\theta \log{\pi_\theta(a|s)}q_{\pi_\theta}(S,A) \Big]
        \end{equation}
    \end{thm}
\section{REINFORCE}
    \subsection{Main Idea}
        \textbf{REINFORCE} is the most basic form of policy gradient method.
        It's the vanilla policy gradient method.
    \subsection{Objective Function}
        REINFORCE uses following objective function;
        \begin{equation}
            J(\pi_\param)=\E_{\tau \sim \pi_\param}[G_0] = \E_{\tau \sim \pi_\param}\left[ \sum_{t=0}^{T-1}\gamma^tr_{t+1} \right].
        \end{equation}
        We use some tricks to compute the gradient of $J(\pi_\param)$
        % TODO: add reference to trick here.
        Then we get
        \begin{equation}
            \nabla_\param J(\pi_\param) = \E_{\tau \sim \pi_\param}\left[ \sum_{t=0}^{T-1}G_t\nabla_\param\log{\pi_\param(a_t|s_t)} \right].
        \end{equation}
        Now we can use this to update the parameter $\param$. Since we cannot compute the exact expectation, we use Monte Carlo sampling
        to approximate $\nabla_\param J(\pi_\param)$ as following:
        \begin{equation}
            \nabla_\param J(\pi_\param) \approx \sum_{t=0}^{T-1}G_t\nabla_\param\log{\pi_\param(a_t|s_t)}
        \end{equation}
    \subsection{Pseudocode}
        \Algorithm{reinforce}
        {REINFORCE}
        {
            \inputitem{$\param$}{initial parameter} \\
            \inputitem{$\gamma$}{discount rate} \\
            \inputitem{$\alpha$}{learning rate}
        }
        {
            \inputitem{$\pi_\param$}{learned policy}
        }
        {
            initialize policy $\pi_\param$\;
            \Repeat{convergence}
            {
                generate episode $\{s_0,a_0,r_1,\dots,s_{T-1},a_{T-1},r_{T}\}$ following $\pi_\param$\;
                \For{$t=0,1,\dots,T-1$}
                {
                    \tcp{compute return-to-go}
                    $G_t=\sum_{k=t+1}^{T} \gamma^{k-t-1}r_k$\;
                    \tcp{policy update}
                    $\param \leftarrow \param + \alpha\nabla_\param\log{\pi_\param(a_t|s_t)}G_t$\;
                }
            }
            \Return{$\pi_\param$}
        }
    \subsection{Implementation Notes}
        \begin{itemize}
            \item When computing return-to-go, we can iterate backward from the terminal time to start, since $G_t \leftarrow \gamma G_{t+1} + r_{t+1}$ where $t=0, 1, 2,\dots,T-1$.
        \end{itemize}
        % TODO: add Implementation details
    \subsection{Pros and Cons}
        Pros:
        \begin{itemize}
            \item No need for model;
            \item Can cover both discrete and continuous actions;
        \end{itemize}
        Cons:
        \begin{itemize}
            \item Large variance;
            \item May converge to suboptimal policy;
        \end{itemize}
    \chapter{Actor-Critic Methods}
    \section{Introduction}
    In this chapter, \textbf{actor-critic} algorithms are introduced. The actor refers to the policy
    update step, and critic refers to the value update step. The actor-critic methods are kind of policy gradient methods,
    with value function method mixed in. Instead of using the full return $G_t$ to compute the loss, which has big variance,
    actor-critic uses value function for the estimate, and this value function is called \textbf{critic}.
    \begin{equation}
        \theta_{t+1} = \theta_t + \alpha \nabla_\theta \log{\pi(a_t|s_t,\theta_t)q_t(s_t,a_t)}
    \end{equation}
    \begin{itemize}
        \item If $q_t(s_t,a_t)$ is estimated by the $G_t$, the corresponding algorithm is the REINFORCE~\ref{alg:reinforce}.
        \item If $q_t(s_t,a_t)$ is estimated by using the TD algorithm, the corresponding algorithms are usually called \textbf{actor-critic}.
    \end{itemize}

    From here, the most simple actor-critic algorithm can be derived, the \textbf{Q actor-critic} algorithm.
    We just use two networks, one for policy and one for value function.
    
    \Algorithm{qac}{QAC}
    {
        \inputitem{$\pi_\theta(a|s)$}{the policy function with initial parameter $\theta_0$} \\
        \inputitem{$q(s, a, w)$}{the value function with initial parameter $w_0$}
    }
    {
        \inputitem{$\pi^*$}{(sub)optimal policy}
    }
    {
        pick an action $a_0$ following $\pi_{\theta_0}(a|s_0)$\;
        \ForEach{time step $t$}
        {
            take $a_t$\;
            observe $r_{t+1}$ and $s_{t+1}$, and then generate $a_{t+1}$ following $\pi_{\theta_t}(a|s_{t+1})$\;
            \tcp{policy update (actor)}
            $\theta_{t+1} \leftarrow \theta_t + \alpha_\theta \nabla_\theta \log{\pi_{\theta_t}(a_t|s_t)}q_{w_t}(s_t,a_t)$\;
            \tcp{value update (critic)}
            $w_{t+1} \leftarrow w_t + \alpha_w[r_{t+1} + \gamma q_{w_t}(s_{t+1},a_{t+1}) - q_{w_t}(s_t,a_t)]\nabla_w q_w(s_t,a_t)$\;
        }
        \Return{$\pi$}
    }
\section{Advantage Actor-Critic (A2C)}
    \subsection{Main Idea}
        The core idea of A2C is to reduce the variance by introducing the baseline.
        The following equation enable us to use the baseline when estimating state values.
        \begin{equation}
            \E_{S \sim \eta, A \sim \pi}[\nabla_\theta \log{\pi_\theta(A|S)}q_\pi(S,A)] = \E_{S \sim \eta, A \sim \pi}[\nabla_\theta\log{\pi_\theta(A|S)}(q_\pi(S,A)-b(S))]
        \end{equation}
        where $b(S)$ is a scalar function of $S$.
        \begin{proof}
            \begin{align}
                \E_{S \sim \eta, A \sim \pi}\Big[ \nabla_\theta \log{\pi_\theta\log{\pi(A|S)}b(S)} \Big]
                    &= \sum_{s \in \mathcal{S}} \eta(s) \sum_{a \in \mathcal{A}(s)}\pi_\theta(a|s)\nabla_\theta \log{\pi_\theta(a|s)}b(s) \\
                    &= \sum_{s \in \mathcal{S}} \eta(s) \sum_{a \in \mathcal{A}(s)}\nabla_\theta \pi_\theta(a|s) b(s) \\
                    &= \sum_{s \in \mathcal{S}} \eta(s)b(s) \sum_{a \in \mathcal{A}(s)} \nabla_\theta \pi_\theta(a|s) \\
                    &= \sum_{s \in \mathcal{S}} \eta(s)b(s) \nabla_\theta \sum_{a \in \mathcal{A}(s)}\pi_\theta(a|s) \\
                    &= \sum_{s \in \mathcal{S}} \eta(s)b(s) \nabla_\theta 1 = 0.
            \end{align}
        \end{proof}

        The baseline reduces the variance of estimation of gradients. In particular, let
        \begin{equation}
            X(S,A) = \nabla_\theta \log{\pi_\theta(a|s)}[q_\pi(S,A)-b(S)]
        \end{equation}
        Then, the true gradient is $\E[X(S,A)]$ and we estimate it with samples $x \sim X$. To better estimate the gradient,
        it would be favorable to make the variance of $X$ small. The theoretically optimal baseline that minimizes $\text{var}(X)$ is
        \begin{equation}
            b^*(s) = \frac{\E_{A \sim \pi}\big[ \|\nabla_\theta \log_\theta(A|s)\|^2 q_\pi(s,A) \big]}
            {\E_{A \sim \pi}\big[ \|\nabla_\theta\log{\pi(A|s)}\|^2 \big]}
        \end{equation}
        proof can be found at \cite{Zhao2025}.
        But it is too complex to be useful in practice. In practice, we use the suboptimal version without the logarithm term, which is actually a value function, intersetingly.
        \begin{equation}
            b(s) = \E_{A \sim \pi}[q_\pi(s,A)] = v_\pi(s)
        \end{equation}

    \subsection{Pseudocode}
        \Algorithm{a2c}
        {Advantage Actor-Critic Algorithm}
        {
            \inputitem{$\pi_\theta$}{a policy function with initial parameter $\theta_0$} \\
            \inputitem{$v_w(s)$}{a value function with initial parameter $w_0$}
        }
        {
            \inputitem{$\pi^*$}{(sub)optimal policy}
        }
        {
            \ForEach{time step $t$} {
                generate and take an action $a_t$ following $\pi_\theta$\;
                observe $r_{t+1}$ and $s_{t+1}$.\;
                \tcp{advantage}
                $\delta_t \leftarrow r_{t+1} + \gamma v_{w_t}(s_{t+1}) - v_{w_t}(s_t)$\;
                \tcp{actor (policy update)}
                $\theta_{t + 1} \leftarrow \theta_t + \alpha_\theta \delta_t \nabla_\theta \log{\pi_{\theta_t}(a_t|s_t)}$\;
                \tcp{critic (value update)}
                $w_{t+1} \leftarrow w_t + \alpha_w \delta_t \nabla_w v_{w_t}(s_t)$\; 
            }
            \Return{$\pi$}
        }
        
        Using the baseline $b(s) = v_\pi(s)$, we get the following update equation:
        \begin{equation}
            \theta_{t+1} = \theta_t + \alpha \nabla_\theta \log{\pi_{\theta_t}(a_t|s_t)}[q_t(s_t,a_t) - v(s_t)]
        \end{equation}

        If $q_t(s_t,a_t)$ and $v_t(s_t)$ were estimated by Monte Carlo, the algorithm is called \textbf{REINFORCE with a baseline}.
        If estimated by TD learining, the algorithm~\ref{alg:a2c} is called \textbf{A2C} or \textbf{TD actor-critic}.
        Note that algorithm~\ref{alg:a2c} estimates $q_t(s_t,a_t) - v_t(s_t)$ with TD error, which is $r_{t+1} + \gamma v_t(s_{t+1} - v_t(s_t))$.
    
    \subsection{Implementation Details}
        We can use n-step update to balance between bias and variance of estimation.
        \begin{equation}
            G_t^{(n)} = r_{t + 1} + \gamma r_{t+2} + \gamma^2 r_{t+3} + \dots + \gamma^n v_t(s_{n+1})
        \end{equation}

        To stop the policy from being almost deterministic too quickly, we can add \textbf{entropy bonus} at the loss function (\cite{pmlr-v48-mniha16}, \cite{WILLIAMS1991}).
        Specifically, the loss function is
        \begin{equation}
            \mathcal{L} = \frac{1}{N}\sum_{i = 0}^{N-1} \Big[-\log{\pi_\theta(a_{t+i}|s_{t+i})}A_{t+i} + c_v(v_\phi(s_{t+i}) - G_{t+i})^2 - c_e \mathcal{H}(\pi_\theta(\cdot|s_{t+i}))\Big]
        \end{equation}
        where 
        \begin{align}
            \mathcal{H}(\pi(\cdot|s)) &= -\int \pi(a|s)\log{\pi(a|s)}da \\
            \mathcal{H} &= -\sum_a \pi(a|s)\log{\pi(a|s)},
        \end{align}
        $A$ and $G$ are advantage and return estimate, respectively.

        The actor and critic can share parameters, when they are using the encoder. The can use different encoder, each of their own, or share one encoder.
    \subsection{Pros and Cons}
    Pros:
    \begin{itemize}
        \item Requires only one networks, opposed to QAC;
        \item Variance is decreased, opposed to the original REINFORCE without baseline;
        \item Appliable to both discrete and continuous actions;
    \end{itemize}
    Cons:
    \begin{itemize}
        \item If the baseline is incorrect, it can be biased;
        \item Sample efficiency is low, since it is an on-policy method;
        \item Vulnerable to big, sudden change of parameters. \textbf{PPO} was designed to resolve this;
    \end{itemize}
    \chapter{Model-based Methods}
    \chapter{Comparisons of Algorithms}
\part{Appendices}
    \chapter{Probability Theory}\label{ch:Prob}
    \section{Introduction}
    This section introduces a theory of measure-theoretic probability, based on \cite{Brzezniak2007}.
\section{Measure-Theoretic Probability}
    \subsection{Basic Definitions}
        \begin{defn}{$\sigma$-field}{sigma_field}
            Let $\Omega$ be a non-empty set. A $\sigma$-field on $\Omega$ is a family of subsets of $\Omega$ such that
            \begin{enumerate}
                \item the empty set $\emptyset \in \mathcal{F}$;
                \item if $A \in \mathcal{F}$ then $\Omega \setminus A \in \mathcal{F}$;
                \item if $A_1, A_2, \dots$ is a sequence of sets in $\mathcal{F}$, then $A_1 \cup A_2 \cup \dots \in \mathcal{F}$ 
            \end{enumerate}
        \end{defn}
        
        \begin{defn}{Probability Measure}{prob_measure}
            Let $\mathcal{F}$ be a $\sigma$-field on $\Omega$. A \textbf{probability measure} $P$ is a function $P: \mathcal{F} \to [0,1]$
            such that
            \begin{enumerate}
                \item $P(\Omega)=1$;
                \item if $A_1, A_2, \dots$ are pairwise disjoint sets ($A_i \cap A_j = \emptyset$ for all $i \neq j$) beloning to $\mathcal{F}$, then
                \begin{equation}
                    P(A_1 \cup A_2 \cup \dots) = P(A_1) + P(A_2) + \dots.
                \end{equation}
            \end{enumerate}
            The triple $(\Omega, \mathcal{F}, P)$ is called a \textbf{probability space}. The sets belonging to $\mathcal{F}$ are called \textbf{events}.
            An event $A$ is said to occur \textbf{almost surely} (a.s.) whenever $P(A)=1$.
        \end{defn}

        \begin{defn}{Random Variable}{rand_var}
            Let $\mathcal{F}$ be a $\sigma$-field on $\Omega$. Then a function $\xi:\Omega \to \mathbb{R}$ is said to be
            $\mathcal{F}$-measurable if the following is satisfied:
            \begin{equation}
                \{\xi \in B\} \in \mathcal{F} \text{ for every Borel set }B \in \mathcal{B}(\mathbb{R}).
            \end{equation}
            where \textbf{Borel sets} refers to the elements of $\mathcal{B}(\mathbb{R})$ which is the smallest $\sigma$-field containing all intervals in $\mathbb{R}$.
            If the triple $(\Omega, \mathcal{F}, P)$ is a probability space, then such a function $\xi$ is called a \textbf{random variable}.
        \end{defn}
        \begin{rem}{Short-hand Notation for Events}{short_hand_notation}
            The set of events such as $\{\xi \in B\}$ represents the set $\{\omega \in \Omega | \xi(\omega) \in B\}$.
        \end{rem}
        
    
    \chapter{Linear Algebra}
    \section{Introduction}
A short useful theorems of linear algebra.
\section{Geshgorin Circle Theorem}
\begin{thm}{Geshgorin Circle Theorem}{Geshgorin}
\end{thm}
\printbibliography
\end{document}